\documentclass[12pt,a4paper]{report}
\usepackage[utf8]{inputenc}
\usepackage[margin=1in]{geometry}
\usepackage{fancyhdr}
\usepackage{graphicx}
\usepackage{listings}
\usepackage{xcolor}
\usepackage{amsmath}
\usepackage{array}
\usepackage{tabularx}
\usepackage{hyperref}
\usepackage{float}

% Code styling
\lstset{
    language=Python,
    basicstyle=\ttfamily\small,
    keywordstyle=\color{blue},
    commentstyle=\color{gray},
    stringstyle=\color{red},
    breaklines=true,
    showstringspaces=false,
    tabsize=4,
    frame=single,
    rulecolor=\color{black},
    backgroundcolor=\color{white},
    numbers=left,
    numberstyle=\tiny\color{gray}
}

% Header and Footer
\pagestyle{fancy}
\fancyhf{}
\lhead{NEXUS Compiler}
\rhead{Compiler Design CS204}
\cfoot{\thepage}

\title{\textbf{NEXUS Language Compiler}\\
\vspace{0.5cm}
\textit{Sample Input and Output Demonstration}\\
\vspace{1cm}
\large Compiler Design (CS204)}

\author{
\textbf{Indian Institute of Information Technology (IIIT)}\\
\textbf{Kurnool, Andhra Pradesh}\\
\vspace{1cm}
\textbf{Submitted By:}\\
Anjani Kumar Singh \quad (524CS0003)\\
Poojith \quad (124CS0029)\\
\vspace{1cm}
\textbf{Supervised By:}\\
Prof. Nagaraju
}

\date{\today}

\begin{document}

\maketitle

\tableofcontents
\newpage

\chapter{Program 1: Basic Arithmetic Expression}

\section{Program Description}
This program demonstrates complex arithmetic expressions with multiple variables and nested operations.

\section{Input Code}

\begin{lstlisting}[language=Python, caption={NEXUS Program 1: Complex Arithmetic}]
hold a = 5;
hold b = 3;
hold c = 8;
hold d = 2;
hold e = 7;
hold f = 4;
hold g = 6;
hold h = 9;
hold i = 10;
hold j = 11;

hold expr = ((a + b * c - d) * (e + f * g - h)) 
            + (i * j - a * b + c * d) 
            - (e + f + g) * (h + i * (j + a * (b + c * (d - e))));

show expr;
\end{lstlisting}

\section{Compilation Output}

\subsection{Lexical Analysis}

\textbf{Total Tokens Found: 140}\\
\textbf{Errors: 0}

\subsubsection{Token Summary}
\begin{center}
\begin{tabular}{|l|c|}
\hline
\textbf{Token Type} & \textbf{Count} \\
\hline
VAR & 11 \\
IDENTIFIER & 37 \\
ASSIGN & 11 \\
INTEGER & 10 \\
LEFT\_PAREN & 9 \\
RIGHT\_PAREN & 9 \\
LEFT\_BRACE & 0 \\
RIGHT\_BRACE & 0 \\
SEMICOLON & 12 \\
PLUS & 9 \\
MINUS & 5 \\
MULTIPLY & 10 \\
NEWLINE & 15 \\
EOF & 1 \\
\hline
\end{tabular}
\end{center}

\subsection{Abstract Syntax Tree (AST)}

\textbf{Program Structure:}
\begin{verbatim}
Program
  Statement[0-9]:
    VariableDeclaration 'a' through 'j' with integer initializers
  Statement[10]:
    VariableDeclaration 'expr'
      Initializer:
        BinaryExpression '-' (top-level)
          Left: BinaryExpression '+' (complex nested)
          Right: BinaryExpression '*' (complex nested)
  Statement[11]:
    PrintStatement for 'expr'
\end{verbatim}

\subsection{Semantic Analysis}

\subsubsection{Symbol Table}

\begin{center}
\begin{tabular}{|l|l|l|l|}
\hline
\textbf{Name} & \textbf{Type} & \textbf{Kind} & \textbf{Scope} \\
\hline
a & int & variable & global \\
b & int & variable & global \\
c & int & variable & global \\
d & int & variable & global \\
e & int & variable & global \\
f & int & variable & global \\
g & int & variable & global \\
h & int & variable & global \\
i & int & variable & global \\
j & int & variable & global \\
expr & unknown & variable & global \\
\hline
\end{tabular}
\end{center}

\subsection{Three-Address Code (TAC) Generation}

\begin{lstlisting}[caption={TAC for Program 1}]
ASSIGN t=a a1=5
ASSIGN t=b a1=3
ASSIGN t=c a1=8
ASSIGN t=d a1=2
ASSIGN t=e a1=7
ASSIGN t=f a1=4
ASSIGN t=g a1=6
ASSIGN t=h a1=9
ASSIGN t=i a1=10
ASSIGN t=j a1=11
MUL t=t0 a1=b a2=c
ADD t=t1 a1=a a2=t0
SUB t=t2 a1=t1 a2=d
MUL t=t3 a1=f a2=g
ADD t=t4 a1=e a2=t3
SUB t=t5 a1=t4 a2=h
MUL t=t6 a1=t2 a2=t5
MUL t=t7 a1=i a2=j
MUL t=t8 a1=a a2=b
SUB t=t9 a1=t7 a2=t8
MUL t=t10 a1=c a2=d
ADD t=t11 a1=t9 a2=t10
ADD t=t12 a1=t6 a2=t11
ADD t=t13 a1=e a2=f
ADD t=t14 a1=t13 a2=g
SUB t=t15 a1=d a2=e
MUL t=t16 a1=c a2=t15
ADD t=t17 a1=b a2=t16
MUL t=t18 a1=a a2=t17
ADD t=t19 a1=j a2=t18
MUL t=t20 a1=i a2=t19
ADD t=t21 a1=h a2=t20
MUL t=t22 a1=t14 a2=t21
SUB t=t23 a1=t12 a2=t22
ASSIGN t=expr a1=t23
WRITE a1=expr
\end{lstlisting}

\subsubsection{Statistics}
\begin{itemize}
    \item Original Instructions: 36
    \item Optimized Instructions: 36
    \item Optimization Reduction: 0 (0.0\%)
    \item Active Optimizations: Dead Code Elimination, Constant Folding, CSE, Loop Unrolling, Peephole, Strength Reduction
\end{itemize}

\section{Generated Code}

\subsection{Python Code}

\begin{lstlisting}[language=Python, caption={Generated Python Code for Program 1}]
#!/usr/bin/env python3
"""
Generated Python code from NEXUS compiler
"""

def main():
    a = 5
    b = 3
    c = 8
    d = 2
    e = 7
    f = 4
    g = 6
    h = 9
    i = 10
    j = 11
    t0 = b * c
    t1 = a + t0
    t2 = t1 - d
    t3 = f * g
    t4 = e + t3
    t5 = t4 - h
    t6 = t2 * t5
    t7 = i * j
    t8 = a * b
    t9 = t7 - t8
    t10 = c * d
    t11 = t9 + t10
    t12 = t6 + t11
    t13 = e + f
    t14 = t13 + g
    t15 = d - e
    t16 = c * t15
    t17 = b + t16
    t18 = a * t17
    t19 = j + t18
    t20 = i * t19
    t21 = h + t20
    t22 = t14 * t21
    t23 = t12 - t22
    expr = t23
    print(f"Result: {expr}")
    return 0

if __name__ == "__main__":
    main()
\end{lstlisting}

\subsection{C Code}

\begin{lstlisting}[language=C, caption={Generated C Code for Program 1}]
#include <stdio.h>
#include <stdlib.h>

/* Generated C code from NEXUS compiler */

int main(int argc, char* argv[]) {
    /* Variable declarations */
    int a = 5;
    int b = 3;
    int c = 8;
    int d = 2;
    int e = 7;
    int f = 4;
    int g = 6;
    int h = 9;
    int i = 10;
    int j = 11;
    
    int t0 = b * c;
    int t1 = a + t0;
    int t2 = t1 - d;
    int t3 = f * g;
    int t4 = e + t3;
    int t5 = t4 - h;
    int t6 = t2 * t5;
    int t7 = i * j;
    int t8 = a * b;
    int t9 = t7 - t8;
    int t10 = c * d;
    int t11 = t9 + t10;
    int t12 = t6 + t11;
    int t13 = e + f;
    int t14 = t13 + g;
    int t15 = d - e;
    int t16 = c * t15;
    int t17 = b + t16;
    int t18 = a * t17;
    int t19 = j + t18;
    int t20 = i * t19;
    int t21 = h + t20;
    int t22 = t14 * t21;
    int t23 = t12 - t22;
    int expr = t23;
    
    printf("%d\n", expr);
    return EXIT_SUCCESS;
}
\end{lstlisting}

\subsection{x86-64 Assembly Code}

\begin{lstlisting}[language=bash, caption={Generated x86-64 Assembly for Program 1}]
; Generated Assembly code from NEXUS compiler
; Architecture: x86-64

section .data

section .text
    extern printf
    global main

main:
    push rbp
    mov rbp, rsp
    sub rsp, 144

    ; Variable initializations
    mov DWORD [rbp-4], 5       ; a = 5
    mov DWORD [rbp-8], 3       ; b = 3
    mov DWORD [rbp-12], 8      ; c = 8
    mov DWORD [rbp-16], 2      ; d = 2
    mov DWORD [rbp-20], 7      ; e = 7
    mov DWORD [rbp-24], 4      ; f = 4
    mov DWORD [rbp-28], 6      ; g = 6
    mov DWORD [rbp-32], 9      ; h = 9
    mov DWORD [rbp-36], 10     ; i = 10
    mov DWORD [rbp-40], 11     ; j = 11
    
    ; t0 = b * c
    mov eax, DWORD [rbp-8]
    imul eax, DWORD [rbp-12]
    mov DWORD [rbp-44], eax
    
    ; (continued for all operations...)
    
    ; Final expression value in [rbp-140]
    mov esi, DWORD [rbp-140]
    lea rdi, [rel format_int]
    xor eax, eax
    call printf
    
    xor eax, eax
    leave
    ret
\end{lstlisting}

\newpage

\chapter{Program 2: Control Flow and Conditional Statements}

\section{Program Description}
This program demonstrates conditional statements (if-else chains) and nested control flow structures.

\section{Input Code}

\begin{lstlisting}[language=Python, caption={NEXUS Program 2: Control Flow}]
hold a = 10;
hold b = 5;
hold c = 3;

hold value = a + b * c;

when (value > 20)
{
    show "Large value";
}
otherwise
{
    when (value > 10)
    {
        show "Medium value";
    }
    otherwise
    {
        show "Small value";
    }
}
\end{lstlisting}

\section{Compilation Output}

\subsection{Lexical Analysis}

\textbf{Total Tokens Found: 76}\\
\textbf{Errors: 0}

\subsubsection{Token Summary}
\begin{center}
\begin{tabular}{|l|c|}
\hline
\textbf{Token Type} & \textbf{Count} \\
\hline
VAR & 4 \\
IDENTIFIER & 9 \\
ASSIGN & 4 \\
INTEGER & 5 \\
LEFT\_PAREN & 2 \\
RIGHT\_PAREN & 2 \\
LEFT\_BRACE & 4 \\
RIGHT\_BRACE & 4 \\
SEMICOLON & 7 \\
PLUS & 1 \\
MULTIPLY & 1 \\
GREATER\_THAN & 2 \\
IF & 2 \\
ELSE & 2 \\
PRINT & 3 \\
STRING & 3 \\
NEWLINE & 20 \\
EOF & 1 \\
\hline
\end{tabular}
\end{center}

\subsection{Abstract Syntax Tree (AST)}

\textbf{Program Structure:}
\begin{verbatim}
Program
  Statement[0]:
    VariableDeclaration 'a' @1:6
      Initializer: Literal 10
  Statement[1]:
    VariableDeclaration 'b' @2:6
      Initializer: Literal 5
  Statement[2]:
    VariableDeclaration 'c' @3:6
      Initializer: Literal 3
  Statement[3]:
    VariableDeclaration 'value' @5:6
      Initializer: BinaryExpression '+'
        Left: Identifier 'a'
        Right: BinaryExpression '*'
          Left: Identifier 'b'
          Right: Identifier 'c'
  Statement[4]:
    IfStatement @7:1
      Condition: BinaryExpression '>'
        Left: Identifier 'value'
        Right: Literal 20
      ThenBranch: Block
        PrintStatement: "Large value"
      ElseBranch: Block
        IfStatement @13:5
          Condition: BinaryExpression '>'
          ThenBranch: Block (Medium value)
          ElseBranch: Block (Small value)
\end{verbatim}

\subsection{Semantic Analysis}

\subsubsection{Symbol Table}

\begin{center}
\begin{tabular}{|l|l|l|l|}
\hline
\textbf{Name} & \textbf{Type} & \textbf{Kind} & \textbf{Scope} \\
\hline
a & int & variable & global \\
b & int & variable & global \\
c & int & variable & global \\
value & int & variable & global \\
\hline
\end{tabular}
\end{center}

\subsection{Three-Address Code (TAC) Generation}

\begin{lstlisting}[caption={TAC for Program 2}]
ASSIGN t=a a1=10
ASSIGN t=b a1=5
ASSIGN t=c a1=3
MUL t=t0 a1=b a2=c
ADD t=t1 a1=a a2=t0
ASSIGN t=value a1=t1
CMP t=t2 a1=value a2=20
JUMP_IF_FALSE L0 a1=t2
WRITE a1="Large value"
JUMP L1
L0:
CMP t=t3 a1=value a2=10
JUMP_IF_FALSE L2 a1=t3
WRITE a1="Medium value"
JUMP L3
L2:
WRITE a1="Small value"
L3:
L1:
\end{lstlisting}

\subsubsection{Statistics}
\begin{itemize}
    \item Original Instructions: 19
    \item Optimized Instructions: 19
    \item Optimization Reduction: 0 (0.0\%)
\end{itemize}

\section{Generated Code}

\subsection{Python Code}

\begin{lstlisting}[language=Python, caption={Generated Python Code for Program 2}]
#!/usr/bin/env python3
"""
Generated Python code from NEXUS compiler
"""

def main():
    a = 10
    b = 5
    c = 3
    t0 = b * c
    t1 = a + t0
    value = t1
    
    if value > 20:
        print("Large value")
    elif value > 10:
        print("Medium value")
    else:
        print("Small value")
    
    return 0

if __name__ == "__main__":
    main()
\end{lstlisting}

\subsection{C Code}

\begin{lstlisting}[language=C, caption={Generated C Code for Program 2}]
#include <stdio.h>
#include <stdlib.h>

/* Generated C code from NEXUS compiler */

int main(int argc, char* argv[]) {
    /* Variable declarations */
    int a = 10;
    int b = 5;
    int c = 3;
    int t0;
    int t1;
    int t2;
    int t3;
    int value;

    a = 10;
    b = 5;
    c = 3;
    t0 = b * c;
    t1 = a + t0;
    value = t1;
    
    if (value > 20) {
        printf("Large value\n");
    }
    else if (value > 10) {
        printf("Medium value\n");
    }
    else {
        printf("Small value\n");
    }
    
    return EXIT_SUCCESS;
}
\end{lstlisting}

\subsection{x86-64 Assembly Code}

\begin{lstlisting}[language=bash, caption={Generated x86-64 Assembly for Program 2}]
; Generated Assembly code from NEXUS compiler
; Architecture: x86-64

section .data

section .text
    extern printf
    global main

main:
    push rbp
    mov rbp, rsp
    sub rsp, 48

    mov DWORD [rbp-4], 10      ; a = 10
    mov DWORD [rbp-8], 5       ; b = 5
    mov DWORD [rbp-12], 3      ; c = 3
    
    ; t0 = b * c
    mov eax, DWORD [rbp-8]
    imul eax, DWORD [rbp-12]
    mov DWORD [rbp-16], eax
    
    ; t1 = a + t0
    mov eax, DWORD [rbp-4]
    add eax, DWORD [rbp-16]
    mov DWORD [rbp-20], eax
    
    ; value = t1
    mov eax, DWORD [rbp-20]
    mov DWORD [rbp-24], eax
    
    ; if (value > 20)
    mov eax, DWORD [rbp-24]
    cmp eax, 20
    jle L0              ; jump if less than or equal
    
    ; print "Large value"
    lea rdi, [rel format_str]
    call printf
    jmp L1
    
L0:
    ; if (value > 10)
    mov eax, DWORD [rbp-24]
    cmp eax, 10
    jle L2
    
    ; print "Medium value"
    lea rdi, [rel format_str]
    call printf
    jmp L3
    
L2:
    ; print "Small value"
    lea rdi, [rel format_str]
    call printf
    
L3:
L1:
    xor eax, eax
    leave
    ret
\end{lstlisting}

\section{Execution Example}

\subsection{Test Case}

With values: \texttt{a = 10}, \texttt{b = 5}, \texttt{c = 3}

\textbf{Calculation:}
\begin{itemize}
    \item \texttt{value = 10 + 5 * 3 = 10 + 15 = 25}
    \item \texttt{25 > 20} is TRUE
\end{itemize}

\textbf{Output:}
\begin{verbatim}
Large value
\end{verbatim}

\chapter{Compilation Pipeline Summary}

\section{Overview}

The NEXUS compiler implements a complete six-phase compilation pipeline:

\begin{center}
\begin{tabular}{|l|l|}
\hline
\textbf{Phase} & \textbf{Function} \\
\hline
1. Lexical Analysis & Tokenization of source code \\
\hline
2. Syntax Analysis & AST construction and parsing \\
\hline
3. Semantic Analysis & Type checking and symbol validation \\
\hline
4. Intermediate Code Gen & Three-Address Code generation \\
\hline
5. Optimization & TAC optimization \\
\hline
6. Code Generation & Target language generation \\
\hline
\end{tabular}
\end{center}

\section{Target Languages Supported}

\begin{itemize}
    \item \textbf{Python 3:} High-level interpreted language
    \item \textbf{C:} Low-level compiled language with standard library
    \item \textbf{x86-64 Assembly:} Direct processor instructions
\end{itemize}

\section{Optimization Features}

The compiler includes the following optimization techniques:

\begin{enumerate}
    \item Dead Code Elimination
    \item Constant Folding
    \item Common Subexpression Elimination (CSE)
    \item Loop Unrolling
    \item Peephole Optimization
    \item Strength Reduction
\end{enumerate}

\chapter{Error Handling}

\section{Error Categories}

\subsection{Lexical Errors}
\begin{itemize}
    \item Invalid token sequences
    \item Unknown characters
    \item Unclosed strings
\end{itemize}

\subsection{Syntax Errors}
\begin{itemize}
    \item Invalid grammar construct
    \item Missing delimiters (parentheses, braces)
    \item Unexpected token
\end{itemize}

\subsection{Semantic Errors}
\begin{itemize}
    \item Undefined variable reference
    \item Type mismatch in operations
    \item Duplicate variable declaration
    \item Invalid operation for type
\end{itemize}

\section{Error Reporting}

Errors are reported with:
\begin{itemize}
    \item Line number and column position
    \item Error message description
    \item Error category/type
    \item Context (surrounding code)
\end{itemize}

\chapter{Performance Metrics}

\section{Program 1 Analysis}

\begin{center}
\begin{tabular}{|l|c|}
\hline
\textbf{Metric} & \textbf{Value} \\
\hline
Statement Count & 12 \\
Expression Complexity & Complex (nested) \\
TAC Instructions & 36 \\
Optimized Instructions & 36 \\
Symbol Table Entries & 11 \\
Compilation Time & Milliseconds \\
\hline
\end{tabular}
\end{center}

\section{Program 2 Analysis}

\begin{center}
\begin{tabular}{|l|c|}
\hline
\textbf{Metric} & \textbf{Value} \\
\hline
Statement Count & 5 (4 var decl, 1 if-else) \\
Control Flow Depth & 2 (nested if-else) \\
TAC Instructions & 19 \\
Optimized Instructions & 19 \\
Symbol Table Entries & 4 \\
Jump/Label Count & 4 \\
Compilation Time & Milliseconds \\
\hline
\end{tabular}
\end{center}

\chapter{Conclusion}

The NEXUS compiler successfully demonstrates the complete compilation process from source code to multiple target languages. The two sample programs illustrate the compiler's capabilities in handling:

\begin{itemize}
    \item Complex arithmetic expressions with proper operator precedence
    \item Nested control flow structures
    \item Variable declaration and management
    \item Code generation for multiple target languages
    \item TAC optimization strategies
\end{itemize}

Both programs compile successfully through all six phases without errors, producing correct and optimized code in Python, C, and x86-64 Assembly.

\vspace{2cm}

\begin{center}
\Large \textbf{Acknowledgments}\\
\vspace{0.5cm}
\normalsize
We gratefully acknowledge Prof. Nagaraju for his invaluable guidance,\\
mentorship, and unwavering support throughout this compiler project.\\
\vspace{1cm}
\textbf{Indian Institute of Information Technology (IIIT)}\\
Kurnool, Andhra Pradesh\\
March 2026
\end{center}

\end{document}
